% =====================================================================
\section{Implementation Schedule and Milestones}
\label{sec:schedule}
% =====================================================================

The schedule covers the six-week capstone period and is ordered by dependency:
no model can be trained before the data are verified, and no fusion experiment is
meaningful before both single-modality baselines work. Training runs are long,
so they are started as early as possible and overlap with development of the
next component. The optional extensions are scheduled only in slack time after
the core system is complete. \Cref{fig:gantt} shows the timeline and
\cref{tab:milestones} the deliverables and completion criteria of each
milestone.

\begin{figure}[H]
\centering
\begin{tikzpicture}[x=1.65cm, y=0.62cm,
  bar/.style={fill=lsblue!70, draw=lsblue, rounded corners=1pt},
  opt/.style={fill=black!15, draw=black!40, rounded corners=1pt, dashed},
  lab/.style={font=\scriptsize\sffamily, anchor=east}]
% grid
\foreach \w in {0,...,6} \draw[black!20] (\w,0.4) -- (\w,-10.4);
\foreach \w in {1,...,6} \node[font=\scriptsize\sffamily\bfseries] at (\w-0.5,1) {Week \w};
% bars: \y, start, end, label
\foreach \y/\s/\e/\l in {
  0/0/1/{M1 Environment \& dataset verification},
  1/0.5/2/{M2 Preprocessing \& alignment},
  2/1.5/2.5/{M3 Audio-only baseline (B1)},
  3/2/3.5/{M4 Visual encoder, B2, B3},
  4/3/4/{M5 Multimodal fusion (P)},
  5/3.5/5/{M6 Training \& tuning},
  6/4.5/5.5/{M7 Evaluation \& ablations},
  7/4.5/6/{M8 Application integration},
  8/5.5/6/{M9 Testing, docs \& demo}} {
  \draw[bar] (\s,-\y-0.35) rectangle (\e,-\y+0.35);
  \node[lab] at (-0.05,-\y) {\l};
}
\draw[opt] (5,-9-0.35) rectangle (6,-9+0.35);
\node[lab] at (-0.05,-9) {Optional extensions (if time)};
\draw[red!60!black, thick] (5,0.5) -- (5,-10.4) node[below, font=\scriptsize\sffamily, text=red!60!black] {core freeze};
\end{tikzpicture}
\caption{Gantt-style implementation timeline. Dashed bar: optional work, started only
after the core system is frozen at the end of week~5.}
\label{fig:gantt}
\end{figure}

{\small
\setlength{\tabcolsep}{4pt}
\begin{longtable}{@{}p{0.8cm}p{3.0cm}p{1.1cm}p{4.4cm}p{5.1cm}@{}}
\caption{Milestones, dependencies, deliverables and completion criteria.}\label{tab:milestones}\\
\toprule
\textbf{ID} & \textbf{Milestone} & \textbf{Needs} & \textbf{Deliverables} & \textbf{Completion criteria} \\
\midrule
\endfirsthead
\multicolumn{5}{@{}l}{\small\itshape \Cref{tab:milestones} continued}\\
\toprule
\textbf{ID} & \textbf{Milestone} & \textbf{Needs} & \textbf{Deliverables} & \textbf{Completion criteria} \\
\midrule
\endhead
\bottomrule
\endfoot
M1 & Environment setup and dataset verification & -- &
Repository skeleton, pinned environment, raw GRID and MUSAN, integrity report, \texttt{splits.csv} &
Every member runs the test suite; integrity report lists usable utterances per speaker; speaker sets disjoint; GPU access confirmed \\
M2 & Preprocessing and data alignment & M1 &
Cached mouth crops, 16\,kHz audio, manifest, noise mixer, fixed noisy test sets, normalisation statistics, data loader &
Crop contact sheets reviewed; noise-mixing tests pass; one batch loads with consistent lengths; train/infer parity test passes \\
M3 & Baseline implementation (audio-only) & M2 &
Training and evaluation scripts, B1 checkpoint &
B1 overfits 8 utterances; full B1 run completes and is evaluated at all five SNRs \\
M4 & Visual and audio encoder development & M2 &
Visual encoder V1, B2 and B3 checkpoints &
Encoder shape tests pass; B2 validation CER falls clearly below that of an untrained model; B3 evaluated at all SNRs \\
M5 & Multimodal fusion & M3, M4 &
Attention fusion, gate, modality dropout, first P checkpoint &
Fusion unit tests pass; no NaN under full-modality dropout; P trains stably \\
M6 & Training and hyperparameter tuning & M5 &
Tuned P configuration; final checkpoints of B1--B3 and P &
Learning-rate and modality-dropout sweeps logged; best configuration chosen on validation data only \\
M7 & Evaluation and ablation studies & M6 &
Per-SNR WER/CER tables and plots, ablations X1--X3, gate analysis, latency measurements &
All core models evaluated on identical test audio; every number traceable to a config and checkpoint \\
M8 & Application integration & M3 (Mode A), M6 (final) &
Gradio Mode A, live Mode B, \texttt{transcribe.py} &
Mode A captions match the evaluation script on test clips; Mode B runs for 5 minutes without error \\
M9 & Testing, documentation and demonstration & M7, M8 &
Final report, README, demonstration rehearsal, frozen release tag &
Fresh clone installs and reproduces one evaluation; demonstration rehearsed with a recorded fallback clip \\
\end{longtable}
}

\paragraph{Prioritisation.} If the schedule slips, work is dropped in this order:
optional extensions; experiments X4--X6; Mode B live captioning (Mode A remains);
multiple seeds. Baselines B1--B3, the proposed model P, the per-SNR evaluation and
ablations X1--X3 are never dropped, because they constitute the core contribution
described in the proposal.

\paragraph{Team organisation.} Work is divided into three pairs \rec: a data pair
(M1--M2, noise and integrity), a model pair (encoders, fusion, training) and an
application pair (inference script, Gradio interface, live mode, tests), with all
members contributing to evaluation and the report. A short weekly review checks
completion criteria and re-plans.
